\documentclass{article}
\usepackage[utf8]{inputenc}
\usepackage[T1]{fontenc}
\usepackage{geometry}
\geometry{
    a4paper,
    left=2.5cm,
    right=2.5cm,
    top=2.5cm,
    bottom=2.5cm
}
\usepackage{graphicx}
\usepackage{float}
\usepackage{hyperref}
\usepackage[english]{babel}
\usepackage{amsmath}
\usepackage{amssymb}
\usepackage{array}
\usepackage{booktabs}
\usepackage{caption}
\usepackage{subcaption}
\usepackage{enumitem}
\usepackage{fancyhdr}
\usepackage{titlesec}
\usepackage{url}
\usepackage{textcomp}
\usepackage{xcolor}
\usepackage{listings}
\usepackage{multirow}
\usepackage{longtable}
\usepackage{tabularx}
\usepackage{times}

% Set up page style
\pagestyle{fancy}
\fancyhf{}
\rfoot{\thepage}
\renewcommand{\headrulewidth}{0pt}

% Title formatting
\titleformat{\chapter}[display]{\normalfont\bfseries}{\chaptertitle\thechapter}{0pt}{\Huge}
\titlespacing*{\chapter}{0pt}{50pt}{40pt}

% Document info
\title{Template Document}
\author{Academic Support Agent}
\date{\today}

\begin{document}
\maketitle

% The following commands are not to be changed as they prepare the initial pages of the document. Also, they set up the environment for the real content.  Do not move them or erase them; they are needed for a proper presentation of your dissertation

\MainTitle
\ThesisCommittee
\Declaration
\Dedication
\Acknowledgements
\Abstract
\contentsAndLists
\startContent

% Here you call the files that you prepared in order to get a nice and concise master thesis (i.e., the chapters).  Just call on the command as many times as necessary. Note that it is the same for either thesis chapters or appendices.

\insertChapter{"Chapter_01"}
\insertChapter{"Chapter_02"}

																	%%	appendices for proper formatting.
\insertChapter{"Appendix_01"}

% These two last commands are not to be changed either. The first one prepares the bibliography and sets its style while the second inserts the author's vita.
\insertBibliography
\insertVita{vitaFile}
\end{document}